% Einheit 2 von 4 – Parameter aus Lagebedingungen (bank/scharen-von-geraden-und-ebenen/e2.jsonl, zusammenbau v0.1)
%% TODO Zeitmarke Einheit 2: Marken-Zeile nicht lesbar
%% TODO Zweigzeile Teil 1: was hier gelernt wird (Satz in Schülersprache aus dem Einheitstitel) – Einheit 2
\einheitenkopf[e2]{Einheit 2 von 4 · Parameter aus Lagebedingungen}
\zweigzeile{Abitur LK}
\setcounter{aufgabe}{11}

%% TODO Titel: Ich-kann-Satz fehlt in der Bank (Platzhalter: Kettenname) – Nr. 12
%% TODO Anweisung: ein Satz über der ersten Teilaufgabe fehlt in der Bank – Nr. 12
\begin{aufgabe}{Parameter aus Lagebedingungen}
%% TODO \rechenplatz{n}: Schrittzahl des Grundfalls fehlt in der Bank (nur bei fester Schreibform, 2.3 b)
\begin{teile}
\teil Wo steht $k$, was verändert es? $E_k: 3x_1 - x_2 + k x_3 = 2$, $k \in \mathbb{R}$ \\ \kreuz{Stützvektor – die Lage ändert sich, die Richtung bleibt} \\ \kreuz{Richtungsvektor – die Richtung ändert sich} \\ \kreuz{Normalenvektor – die Stellung der Ebene ändert sich} \\ \kreuz{rechte Seite – die Ebenen werden parallel verschoben}
\teil Für welches $k$ ist $E_k: k x_1 + 2x_2 - x_3 = 3$ parallel zu einer Geraden mit Richtungsvektor $(2 | 1 | 4)$? $k = $ \leerfeld
\teil Für welches $k$ steht eine Gerade mit Richtungsvektor $(4 | k | -2)$ senkrecht auf $E: 2x_1 + x_2 - x_3 = 4$? $k = $ \leerfeld
\teil Gibt es ein $k$, für das $g: \vec{x} = (1 | 2 | 1) + r \cdot (1 | -1 | 0)$ in $E_k: k x_1 + x_2 + x_3 = 4$ liegt?
\teil Zeige, dass die Kante von $P(2 | 1 | 3)$ nach $Q(4 | 0 | 3)$ für einen Wert von $k$ in $E_k: x_1 + 2x_2 + x_3 = k$ liegt.
\teil Ein Laserstrahl startet in $R_k(4 | 2k - 3 | -2)$ in Richtung $(-2 | 1 | 1)$ und trifft genau den Mittelpunkt $M$ der Strecke von $A(-1 | 3 | 1)$ nach $B(1 | 7 | -1)$. Bestimme $k$. \hfill (Abitur 2018 LK) \\ $k = $ \leerfeld
\teil Stelle eine Koordinatengleichung der Ebene durch $A(2 | 1 | 0)$, $B(0 | 3 | 1)$, $C(1 | 1 | 2)$ auf und bestimme $k$ so, dass $P_k(k | 3 | 2 - k)$ in ihr liegt. $k = $ \leerfeld
\teil Zwei Seiten eines Quadrats liegen auf Geraden der Schar $g_k$ mit dem Richtungsvektor $(2 | 1 | 2)$; $Q(1 | 0 | 2)$ und $R(4 | 3 | 2)$ sind Ecken des Quadrats. Zeige, dass $Q$ und $R$ keine benachbarten Ecken sind. \hfill (Abitur 2024 LK)
\end{teile}
\end{aufgabe}

%% TODO Titel: Ich-kann-Satz fehlt in der Bank (Platzhalter: Kettenname) – Nr. 13
%% TODO Anweisung: ein Satz über der ersten Teilaufgabe fehlt in der Bank – Nr. 13
\begin{aufgabe}{Parameter aus Lagebedingungen – Fehler finden, Begründen}
\begin{teile}
\teil Finde den Fehler. Für welches $k$ ist $E_k: 2k x_1 + x_2 + (k - 3) x_3 = 4$ parallel zu $g: \vec{x} = (1 | 2 | 1) + r \cdot (1 | 1 | 1)$? \rechnung{2k + 2 + (k - 3) &= 4 \\ 3k - 1 &= 4 \\ k &= \tfrac{5}{3}}
\teil Warum verlangt „$E_k$ parallel zu $g$“ das Skalarprodukt null, „$g$ senkrecht zu $E_k$“ aber die Kollinearität?
\end{teile}
\end{aufgabe}
